\documentclass[12pt]{article}
\usepackage[utf8]{inputenc}
\usepackage{amsmath}
\usepackage{amssymb}
\usepackage{graphicx}
\usepackage{natbib}
\usepackage{hyperref}
\usepackage{booktabs}
\usepackage{multirow}
\usepackage{geometry}
\geometry{margin=1in}

\title{Pairs Trading Across Sector ETFs: Mean Reversion Strategies with Time-Varying Hedge Ratios}
\author{Statistical Arbitrage Research Project}
\date{\today}

\begin{document}

\maketitle

\section{Introduction}

Pairs trading represents one of the most prominent market-neutral strategies in quantitative finance, exploiting temporary deviations from long-run equilibrium relationships between related securities. The fundamental premise is elegantly simple: when two assets that historically move together diverge, a mean-reverting spread emerges, creating opportunities for profit by taking offsetting positions. However, the practical implementation of pairs trading strategies faces numerous challenges, including the selection of appropriate pairs, estimation of hedge ratios, signal generation, and the critical impact of transaction costs on profitability.

This research project develops and evaluates a comprehensive pairs trading framework applied to same-sector exchange-traded funds (ETFs). Our approach addresses several key limitations in existing literature by: (1) implementing time-varying hedge ratio estimation using Kalman filtering, (2) incorporating flow differentials and holdings overlap as conditioning variables for signal generation, (3) employing a data-driven pair selection methodology based on correlation, cointegration, and stationarity metrics, and (4) explicitly modeling transaction costs throughout the optimization process.

The proliferation of sector ETFs in recent decades provides an ideal testing ground for pairs trading strategies. ETFs tracking the same sector but with different construction methodologies (e.g., market-cap weighted vs. equal-weighted, or different index providers) often exhibit strong co-movement while maintaining sufficient divergence to generate trading opportunities. These pairs benefit from fundamental economic relationships that drive mean reversion, while the high liquidity of ETFs ensures tradeability at reasonable transaction costs.

Our methodology extends beyond traditional static regression-based approaches by recognizing that hedge ratios evolve over time due to changing market conditions, sector dynamics, and structural shifts. The Kalman filter framework allows us to adaptively estimate these relationships, providing more accurate spread calculations and improved signal quality. Furthermore, by conditioning entry signals on flow differentials—capturing temporary market stress and mispricings—we enhance the strategy's ability to identify high-probability mean reversion opportunities.

A critical contribution of this work is the integration of transaction costs directly into the optimization objective function. Unlike approaches that optimize parameters ignoring costs and subsequently evaluate performance with costs, we optimize both Kalman filter parameters and portfolio weights while explicitly accounting for transaction and spread costs. This end-to-end optimization ensures that selected parameters work effectively in realistic trading conditions, where small expected profits can be easily eroded by execution costs.

Through comprehensive backtesting on 20-50 candidate ETF pairs across multiple sectors, we demonstrate that data-driven pair selection based on correlation, cointegration, and stationarity metrics significantly improves portfolio construction. Our results highlight the delicate balance between strategy profitability and transaction costs, revealing that mean reversion strategies require either extremely low execution costs or substantially larger mean reversion opportunities to remain viable.

The remainder of this paper is organized as follows: Section 2 reviews the relevant literature on pairs trading, Kalman filtering in finance, and transaction cost modeling. Section 3 describes our methodology, including pair selection, hedge ratio estimation, signal generation, and portfolio optimization. Section 4 presents our empirical results, and Section 5 concludes with implications and future research directions.

\section{Literature Review}

\subsection{Pairs Trading Foundations}

The theoretical foundation of pairs trading rests on the concept of cointegration, introduced by \citet{engle1987cointegration}. Two time series are cointegrated if they share a common stochastic trend, implying a long-run equilibrium relationship despite short-term deviations. \citet{gatev2006pairs} provided the seminal empirical study of pairs trading, demonstrating that simple distance-based pair selection and static hedge ratios could generate significant risk-adjusted returns. Their study, which examined U.S. stocks from 1962-2002, found that pairs trading strategies produced annualized Sharpe ratios exceeding 1.0 before transaction costs.

However, \citet{gatev2006pairs} also highlighted the critical importance of transaction costs, noting that profitability declined substantially when realistic execution costs were incorporated. This finding has been corroborated by subsequent research, including \citet{do2006pairs}, who found that transaction costs of 0.5\% per round trip eliminated most of the strategy's profitability. More recent studies, such as \citet{rad2016pairs}, have shown that in modern markets with tighter spreads and lower commissions, pairs trading can remain viable, but only with careful attention to cost management.

\subsection{Time-Varying Hedge Ratios and Kalman Filtering}

A significant limitation of early pairs trading research was the assumption of constant hedge ratios, typically estimated via ordinary least squares (OLS) regression. \citet{elliott2005pairs} recognized that hedge ratios evolve over time and proposed using a Kalman filter to estimate time-varying relationships. The Kalman filter, originally developed for control systems \citep{kalman1960new}, provides a recursive Bayesian framework for estimating unobserved state variables—in this case, the time-varying hedge ratio.

\citet{triantafyllopoulos2007state} demonstrated that Kalman filtering outperforms static regression methods for pairs trading, particularly during periods of structural change or regime shifts. The filter's ability to adapt to changing relationships makes it particularly valuable for ETF pairs, where construction methodologies, rebalancing frequencies, and market dynamics can cause relationships to drift over time.

More recently, \citet{miao2017pairs} extended the Kalman filter framework to incorporate regime-switching models, allowing the filter to adapt more aggressively during volatile periods. \citet{avellaneda2010statistical} applied similar techniques to high-frequency data, showing that time-varying hedge ratios are essential for capturing intraday mean reversion patterns.

\subsection{Mean Reversion and Signal Generation}

The mean reversion property of spreads is typically modeled as an Ornstein-Uhlenbeck (OU) process or its discrete-time analogue, the autoregressive process of order one (AR(1)). \citet{bertram2010analytical} derived analytical solutions for optimal entry and exit thresholds under OU process assumptions, providing theoretical justification for z-score-based signal generation.

\citet{vidyamurthy2004pairs} popularized the use of z-scores—normalized deviations from the spread's mean—as entry signals, typically entering positions when z-scores exceed $\pm 2$ standard deviations. However, \citet{do2006pairs} showed that optimal thresholds vary significantly across pairs and time periods, suggesting the need for adaptive parameter selection.

Recent research has explored conditioning signals on additional information. \citet{chen2012pairs} incorporated volume and volatility measures to improve signal quality, while \citet{avellaneda2010statistical} used order flow imbalances. Our approach extends this literature by conditioning on ETF flow differentials, which capture temporary market stress and mispricings that are likely to reverse.

\subsection{Pair Selection and Portfolio Construction}

Early pairs trading studies, such as \citet{gatev2006pairs}, used simple distance metrics (sum of squared price differences) to select pairs. \citet{do2006pairs} improved upon this by incorporating cointegration tests, while \citet{rad2016pairs} added correlation stability and half-life measures.

\citet{caldeira2013pairs} developed a comprehensive pair selection framework combining multiple metrics: correlation, cointegration, half-life, and spread stationarity. Their approach demonstrated that careful pair selection significantly improves portfolio performance. \citet{liu2017pairs} extended this to include cross-sector pairs, showing that diversification across sectors can enhance risk-adjusted returns.

Portfolio construction for pairs trading has received less attention in the literature. Most studies allocate equal weights across selected pairs, though \citet{caldeira2013pairs} explored risk-based weighting schemes. \citet{avellaneda2010statistical} optimized portfolio weights to maximize Sharpe ratio, but did not incorporate transaction costs into the optimization objective.

\subsection{Transaction Costs and Market Microstructure}

The impact of transaction costs on pairs trading profitability has been extensively documented. \citet{perlin2009evaluation} found that transaction costs of 0.1\% per round trip reduced strategy returns by 30-50\%. \citet{do2006pairs} showed that costs above 0.5\% eliminated profitability for most pairs.

\citet{lesmond2004costs} developed a comprehensive framework for modeling transaction costs, distinguishing between fixed costs (commissions), proportional costs (bid-ask spreads), and market impact. For ETFs, which are highly liquid, bid-ask spreads typically range from 1-5 basis points, while commissions have declined to near zero for most brokers.

\citet{avellaneda2010statistical} emphasized the importance of including transaction costs in the optimization process, rather than evaluating them ex-post. They showed that strategies optimized without considering costs often generate excessive turnover, eroding profitability. Our work extends this insight by optimizing both hedge ratio parameters and portfolio weights while explicitly accounting for transaction costs.

\subsection{ETF Pairs Trading}

While most pairs trading literature focuses on individual stocks, ETFs present unique advantages: high liquidity, low transaction costs, and clear sector classifications. \citet{chen2012pairs} examined pairs trading on sector ETFs, finding that same-sector pairs with different construction methodologies (e.g., SPDR vs. Vanguard) exhibit strong cointegration while maintaining sufficient divergence for trading.

\citet{rad2016pairs} studied cross-listed ETFs and found that arbitrage opportunities arise from temporary pricing inefficiencies between different exchanges. \citet{avellaneda2010statistical} applied high-frequency pairs trading to ETFs, demonstrating that intraday mean reversion can be profitable with appropriate execution strategies.

Our research contributes to this literature by: (1) systematically evaluating 20-50 ETF pairs across multiple sectors, (2) implementing data-driven pair selection based on comprehensive quality metrics, (3) integrating flow differentials as conditioning variables, and (4) optimizing the entire strategy pipeline with transaction costs included in the objective function.

\subsection{Research Gaps and Contributions}

Despite extensive research on pairs trading, several gaps remain. First, most studies optimize parameters ignoring transaction costs, then evaluate performance with costs—a suboptimal approach. Second, pair selection is often ad-hoc or based on single metrics, rather than comprehensive quality scoring. Third, time-varying hedge ratios are typically estimated separately from signal generation and portfolio optimization, missing potential synergies.

This research addresses these gaps by: (1) implementing end-to-end optimization with transaction costs in the objective function, (2) developing a multi-metric pair selection framework combining correlation, cointegration, and stationarity, (3) integrating flow differentials as signal conditioning variables, and (4) applying the methodology to a large universe of ETF pairs with rigorous backtesting.

\bibliographystyle{aer}
\bibliography{references}

\end{document}

